\section*{Hoofdstuk \ref{ch:b2:retrosynthesis} --- Retrosynthese en strategie in meerdere stappen}

\begin{solution}{exo:b2:retrosynthesis:1}
\begin{itemize}
  \item 2-Fenylpropaan-2-ol: propanon + \ce{PhMgBr}, of acetofenon + \ce{CH3MgBr}.
  \item Pentaan-3-ol: propanal + \ce{CH3CH2MgBr}.
  \item 1-Cyclohexylethanol: ethanal + cyclohexylmagnesiumbromide, of cyclohexaancarbaldehyde +
        \ce{CH3MgBr}.
  \item 2-Methylbutaan-2-ol: propanon + \ce{CH3CH2MgBr}, of butanon + \ce{CH3MgBr}.
  \item 1-Fenylethanol: benzaldehyde + \ce{CH3MgBr}, of ethanal + \ce{PhMgBr}.
\end{itemize}
\end{solution}

\begin{solution}{exo:b2:retrosynthesis:2}
Acceptorsynthon $\mathrm{Ph\overset{+}{C}H{-}OH}$, equivalent benzaldehyde; donorsynthon \ce{CH3-},
equivalent methylmagnesiumbromide (of methyllithium). (De andere snede, met \ce{Ph-} als donor, gebruikt
\ce{PhMgBr} en ethanal.)
\end{solution}

\begin{solution}{exo:b2:retrosynthesis:3}
$0.90 \times 0.85 \times 0.80 \times 0.95 \times 0.70 = 0.407$: ongeveer 41~\%.
\end{solution}

\begin{solution}{exo:b2:retrosynthesis:4}
\ce{LiAlH4}, nodig om de ester tot de alcohol te reduceren, zou ook het keton reduceren. Bescherm het
keton als zijn cyclische acetaal (ethaan-1,2-diol, zure katalysator, water verwijderd); reduceer de ester met
\ce{LiAlH4}; verwijder het acetaal met verdund waterig zuur: 5-hydroxypentaan-2-on.
\end{solution}

\begin{solution}{exo:b2:retrosynthesis:5}
FGI: butaan-1,3-diol $\Rightarrow$ 3-hydroxybutanal (reductie van het aldehyde, \ce{NaBH4}). \omterm{def:b2:retrosynthesis:disconnection}{Disconnectie} met twee
groepen van dit $\beta$-hydroxyaldehyde aan de binding $\alpha$--$\beta$: twee moleculen ethanal
(\omterm{def:b2:enolates-aldol:aldol}{aldolreactie}). Vooruit: ethanal, verdunde base, koud; daarna \ce{NaBH4}.
\end{solution}

\begin{solution}{exo:b2:retrosynthesis:6}
De cyclohexeenring wordt doorgeknipt aan de twee bindingen \ce{C-C} die in een \omterm{def:b2:frontier-orbitals:diels-alder}{diels-alderreactie} gevormd worden, de
bindingen die allylisch staan ten opzichte van de dubbele binding in de ring, aan de kant van de substituent: buta-1,3-dieen en but-3-een-2-on, met de
acetylgroep op het \omterm{def:b2:frontier-orbitals:diels-alder}{dienofiel}.
\end{solution}

\begin{solution}{exo:b2:retrosynthesis:7}
Nummering C1 (\ce{C=O}), C2=C3, dan C4 met de twee methylgroepen
(\cref{met:b2:conjugate-additions:robinson-disconnection}): de aldolsnede C2--C3 laat een aldehyde
op C3 en een methylketon achter; de michaelsnede C4--C5 laat 2-methylpropanal (C3 zijn carbonylgroep, C4 zijn
$\alpha$-koolstofatoom) en but-3-een-2-on achter.
\end{solution}

\begin{solution}{exo:b2:retrosynthesis:8}
Eén stap: een \omterm{def:b2:aromatic-substitution:friedel-crafts}{friedel-craftsacylering} van benzeen met butanoylchloride en \ce{AlCl3} (geen omlegging,
enkelvoudige acylering). Twee stappen: benzaldehyde met propylmagnesiumbromide, daarna oxidatie van de secundaire
alcohol. De friedel-craftsroute is korter en haar reagentia goedkoper, maar ze gebruikt een volledig equivalent
aluminiumchloride, dat in het waterige afval terechtkomt.
\end{solution}

\begin{solution}{exo:b2:retrosynthesis:9}
Pyridiniumchloorchromaat in dichloormethaan, de swernoxidatie of het Dess--Martin-perjodinaan
geven hexanal; aangezuurd waterig dichromaat of permanganaat geeft hexaanzuur, omdat het
aldehyde in water zijn hydraat vormt, dat opnieuw geoxideerd wordt (\cref{prop:b2:retrosynthesis:mild-oxidation}).
\end{solution}

\begin{solution}{exo:b2:retrosynthesis:10}
Het amine reageert als eerste, onbeschermd, als de alcohol beschermd is: bescherm de alcohol als silylether
(die overleeft de acylering); acyleer het amine; verwijder de silylgroep met fluoride; acyleer de
alcohol. Als het amine vrij moest blijven terwijl de alcohol reageert, zou het beschermd worden als
\textit{tert}-butoxycarbonylcarbamaat (verwijderd door zuur), orthogonaal ten opzichte van de silylether (verwijderd door
fluoride): elk van beide kan vrijgemaakt worden zonder het andere te raken.
\end{solution}

\begin{solution}{exo:b2:retrosynthesis:11}
\ce{CH3COCH2CH2COCH3} is een 1,4-dicarbonylverbinding: knip de centrale binding \ce{C-C} door. De donor is het \omterm{def:b2:enolates-aldol:enolate}{enolaat} van
ethyl-3-oxobutanoaat (natriumethoxide); de acceptor, een $\alpha$-koolstofatoom van een keton, wordt geleverd door
1-broompropaan-2-on, \ce{BrCH2COCH3} (bimoleculaire substitutie van bromide). Daarna verwijderen waterige hydrolyse van
de ester en verhitten \ce{CO2}: hexaan-2,5-dion.
\end{solution}

\begin{solution}{exo:b2:retrosynthesis:12}
Retrosynthese: het methylketon met een \ce{CH2} ernaast is het product van de acetoazijnestersynthese;
knip de binding tussen C3 en C4 door: het \omterm{def:b2:enolates-aldol:enolate}{enolaat} van ethyl-3-oxobutanoaat en het allylische halogenide
1-broom-3-methylbut-2-een, \ce{(CH3)2C=CHCH2Br}. Vooruit: natriumethoxide in ethanol, daarna het halogenide;
daarna waterige base, aanzuren en verhitten, wat de ester hydrolyseert en \ce{CO2} verwijdert
(\cref{prop:b2:enolates-aldol:decarboxylation}): 6-methylhept-5-een-2-on.
\end{solution}

\begin{solution}{pb:b2:retrosynthesis:1}
\textbf{1.} \ce{HO-C6H4-CH2CH2COCH3}, met de \ce{OH} para ten opzichte van de keten: een fenol en een keton.
\textbf{2.} Hydrogenering van een \ce{C=C}: \ce{ArCH2CH2CO} $\Rightarrow$ \ce{ArCH=CHCO}, een
$\alpha,\beta$-onverzadigd keton.
\textbf{3.} (\textit{E})-4-(4-Hydroxyfenyl)but-3-een-2-on, \ce{HOC6H4CH=CHCOCH3}.
\textbf{4.} Een $\alpha,\beta$-onverzadigd carbonyl: een \omterm{def:b2:enolates-aldol:aldol}{aldolcondensatie}, doorgeknipt aan de \ce{C=C}.
\textbf{5.} 4-Hydroxybenzaldehyde en propanon.
\textbf{6.} Het aldehyde heeft geen $\alpha$-waterstofatoom en kan alleen het elektrofiel zijn; propanon, in
overmaat, is de enige bron van \omterm{def:b2:enolates-aldol:enolate}{enolaat} en reageert eerder met het reactievere aldehyde dan met zichzelf.
\textbf{7.} (1) 4-hydroxybenzaldehyde, overmaat propanon, waterig natriumhydroxide; aanzuren. (2)
\ce{H2} over palladium op koolstof, bij kamertemperatuur, tot één equivalent opgenomen is.
\textbf{8.} De groep \ce{CHO} stabiliseert het fenolaat, dus het fenol is minstens even zuur als fenol
($\mathrm pK_a$ 9.99). Bij $\mathrm{pH} > 13$ is $[\ce{ArO-}]/[\ce{ArOH}] > 10^{13-10} = 10^3$: het
fenolaat.
\textbf{9.} De \ce{O-} duwt elektronendichtheid via de ring naar het carbonylkoolstofatoom (een
mesomere donor in parapositie): het aldehyde is minder elektrofiel, en de condensatie trager.
\textbf{10.} Een benzylether: de hydrogenering over palladium die de \ce{C=C} reduceert, splitst ook
de benzylether (hydrogenolyse), en maakt het fenol in dezelfde stap vrij.
\textbf{11.} Drie: benzylering, \omterm{def:b2:enolates-aldol:aldol}{aldolcondensatie}, hydrogenering met ontscherming.
\textbf{12.} Nee: onder milde omstandigheden (kamertemperatuur, palladium) wordt de \omterm{def:b2:huckel:aromatic}{aromatische} ring niet gereduceerd.
\textbf{13.} Niet beschermen: het fenolaat vertraagt de condensatie alleen, wat een langere tijd of opwarmen
compenseert; één bespaarde stap is meer waard dan de winst in snelheid.
\textbf{14.} 4-Joodfenol en but-3-een-2-on.
\textbf{15.} \ce{HOC6H4I + CH2=CHCOCH3 + Et3N -> HOC6H4CH=CHCOCH3 + Et3NH+ + I-}, palladiumkatalysator.
\textbf{16.} Op de eindstandige \ce{CH2}, het $\beta$-koolstofatoom: de arylgroep inserteert aan het minst gehinderde
uiteinde, en $\beta$-hydride-eliminatie geeft daarna het geconjugeerde (\textit{E})-enon.
\textbf{17.} \omterm{def:b2:catalytic-cycles:oxidative-addition}{Oxidatieve additie} van het aryljodide aan \ce{Pd(0)}, coördinatie en insertie van het
alkeen, $\beta$-hydride-eliminatie, en terugvorming van \ce{Pd(0)} door de base, die \ce{HI} opneemt
(\cref{ch:b2:catalytic-cycles}).
\textbf{18.} $162.188/(220.005 + 70.091 + 101.193) = 162.188/391.289 \approx 41.4~\%$.
\textbf{19.} \ce{C7H6O2 + C3H6O -> C10H10O2 + H2O}: $162.188/(122.123 + 58.080) = 162.188/180.203
\approx 90.0~\%$.
\textbf{20.} 100~\%: elk atoom van \ce{H2} en van het enon komt in het product terecht.
\textbf{21.} $164.204/(122.123 + 58.080 + 2.016) = 164.204/182.219 \approx 90.1~\%$.
\textbf{22.} De \ce{C=C}: geconjugeerd en ongehinderd, bindt ze aan palladium en wordt ze daar veel
sneller gehydrogeneerd dan de \ce{C=O} van het keton; stoppen na één equivalent waterstof spaart het keton.
\textbf{23.} Aldolroute: goedkoop aldehyde, propanon en natriumhydroxide, water als bijproduct.
Heckroute: een aryljodide, een palladiumkatalysator, en but-3-een-2-on, zeer giftig.
\textbf{24.} De aldolroute: twee stappen, goedkope en veiligere reagentia, atoomeconomie boven 90~\%.
\textbf{25.} Het vermenigvuldigt het met 0.90: $0.782 \times 0.90 \approx 70~\%$.
\textbf{26.} $0.85 \times 0.92 = \boldsymbol{\approx 78~\%}$ in totaal.
\end{solution}
